\documentclass[10pt,a4paper]{amsart}
\usepackage[margin=19mm]{geometry}
\usepackage{amsmath,amssymb,mathtools}
\usepackage[scr=rsfs,cal=euler]{mathalfa}
\usepackage{xfrac}
\usepackage[unicode,hidelinks,bookmarks=false]{hyperref}
\newcommand{\ie}{i.e.\ }
\newcommand{\cf}{cf.\ }
\newcommand{\resp}{respectively\ }
\newcommand{\Subsection}{\subsection}
\providecommand{\nobreakdash}{\nobreak\hbox{-}}
\setlength{\emergencystretch}{2em}
\multlinegap0pt
\usepackage{bm}
\usepackage{bbm}
\usepackage{dsfont}
\usepackage{xspace}
\usepackage{cancel}
\usepackage{mathtools}
\usepackage{amsthm}
\makeatletter
\@ifundefined{theorem}{\newtheorem{theorem}{Theorem}}{}
\@ifundefined{prop}{\newtheorem{prop}[theorem]{Proposition}}{}
\makeatother
\newcommand{\C}{\mathbb{C}}
\renewcommand{\emptyset}{\varnothing}
\title[Relative Kazhdan Lusztig isomorphism for $GL\_{2n}/Sp\_{2n}$]{Relative Kazhdan Lusztig isomorphism for $GL\_{2n}/Sp\_{2n}$}
\author{Guy Shtotland}
\date{Source: arXiv:2601.22846v1 (CC BY 4.0)}
\begin{document}
\maketitle

\noindent\emph{Source and licence.} Relative Kazhdan Lusztig isomorphism for $GL_{2n}/Sp_{2n}$. Version arXiv:2601.22846v1 (CC BY 4.0), distributed under the Creative Commons Attribution 4.0 International licence, as recorded in the arXiv licence metadata for this exact version and shown on the abstract page https://arxiv.org/abs/2601.22846v1. Full rights notice: licenses/arxiv-2601.22846-CC-BY-4.0.txt.\par\smallskip

\noindent\textit{Editorial scope.} Counted unit: the Prop whose source label is \texttt{None} in the article (source0.tex lines 814--820, source label \texttt{None}), delivered together with the article's own complete original proof of it (source0.tex lines 822--857, one \texttt{proof} environment of 36 lines). No mathematics is written, repaired or re-derived: statement and proof are the article's own text line for line. Required context retained uncounted: Ber\_presentation (lines 773--779); action (lines 1902--1915). Every definition, display and label of those lines is retained unchanged. Omitted from this package and not redistributed: 1--772, 780--813, 821--821, 858--1901, 1916--2008. Nothing inside a retained excerpt is omitted or altered: every retained body is delivered exactly as the article's lines are written, so no omission receipt of any kind accompanies this package. Citations: the retained excerpts carry no citation command at all, so no bibliography entry of the article is transcribed and no reference is redistributed. Reference adaptation: none was needed and none was made; every reference inside the delivered unit resolves inside it, so no unresolved reference or citation is produced. Printed slips kept literal: no printed word, symbol, hypothesis or number of the counted statement or of its proof was altered. Layout and class: the article's own class is \texttt{article} with packages including babel, graphicx, framed, ulem, amsmath, amsthm, inputenc, amssymb, amsfonts, tikz-cd, enumerate, tikz, geometry, enumitem; the wrapper is the lane's pinned \texttt{amsart} editorial wrapper at 10pt on A4, which restates the theorem-like environments the retained text needs and adds the article's own macro definitions that the retained text needs (\texttt{\textbackslash C}, \texttt{\textbackslash emptyset}), with extra wrapper packages (bm, bbm, dsfont, xspace, cancel, mathtools). The delivered numbering is the unit's own. Assets: no retained span contains a figure environment, a graphics-inclusion command, a PDF-page inclusion, a table or a TikZ picture; no image file is copied into this package, the package declares no asset dependency and no image file is redistributed with it. Licence: version arXiv:2601.22846v1 of this article is distributed under the Creative Commons Attribution 4.0 International licence, as recorded in the arXiv licence metadata for this exact version and shown on the abstract page https://arxiv.org/abs/2601.22846v1. A full rights notice is delivered at licenses/arxiv-2601.22846-CC-BY-4.0.txt. Characters: every retained excerpt is pure ASCII, so the delivered engine, XeLaTeX with its default Latin Modern text font, reports no missing character.
\begin{theorem}\label{Ber_presentation}
    Let $H_f$ be the finite Hecke algebra, it is the sub algebra of $H(G,I)$ generated by $T_w$ for $w\in W$ in the finite Weyl group. 

    $H(G,I)$ is isomorphic to the algebra generated by $H_f$ and $\Theta$ with the following relations.

    For $s=s_\alpha\in W$ a simple reflection and $\lambda\in X_*(T)$ we have $$T_s\theta_{s(\lambda)}-\theta_\lambda T_s=(1-q)\cdot \frac{\theta_\lambda-\theta_{s(\lambda)}}{1-\theta_{-\alpha^\vee}}=(1-q)\sum^{\alpha(\lambda)-1}_{i=0} \theta_\lambda\theta^i_{-\alpha^\vee}$$
\end{theorem}

\begin{prop}\label{action}
For an element $w\in W_{aff}$ exactly one of the following three statements holds:
\begin{enumerate}
    \item There exists $t_n\rightarrow\infty$ such that $T_s1_{Iwt_nH}=1_{Iswt_nH}$
    \item There exists $t_n\rightarrow\infty$ such that $T_s1_{Iwt_nH}=q1_{Iswt_nH}+(q-1)1_{Iwt_nH}$
    \item There exists $t_n\rightarrow\infty$ such that $T_s1_{Iwt_nH}=q1_{Iswt_nH}$
\end{enumerate}
And we have respectively
\begin{enumerate}
    \item $T_s1_{IwH_P}=1_{IswH_P}$
    \item $T_s1_{IwH_P}=q1_{IswH_P}+(q-1)1_{IwH_P}$
    \item $T_s1_{IwH_P}=q1_{IswH_P}$
\end{enumerate}
\end{prop}

    \begin{prop}
    The annihilator of $f_\emptyset\in S(X_\emptyset)^I$ is generated by the elements of the form:
    \begin{enumerate}
        \item $T_{s_{2i-1}}-q$.
        \item $\theta_{\lambda}-q^\frac{l(\lambda^-)-l(\lambda^+)}{2}$ for $\lambda\in L_{max}$.
    \end{enumerate}
    \end{prop}

    \begin{proof}
        We begin with checking that these elements indeed annihilate $f_\emptyset$.

        Let $s=s_{2i-1}$, by Proposition \ref{action} to see that $T_sf_\emptyset=qf_\emptyset$ we need to find a sequence $t_m\in T$, $t_m\rightarrow\infty$ such that $Ist_mSp_{2n}=It_mSp_{2n}$. This holds for any $t\in T$.

        Now, let $\lambda\in L_{max}$, we need to show that $\theta_{\lambda}f_\emptyset= q^\frac{l(\lambda^-)-l(\lambda^+)}{2}f_\emptyset$. We claim that for every $\lambda\in X_*(T)$, $\theta_{\lambda}f_\emptyset=\theta_{\lambda}1_{IS_\emptyset}= q^\frac{l(\lambda^-)-l(\lambda^+)}{2}1_{I\lambda(u)S_\emptyset}$. 
        
        In order to prove this, it is enough to show that if $\lambda\in X_*(T)$ is dominant and $t'\in T$ is arbitrary then $1_{I\lambda(u)I}1_{It'S_\emptyset}=1_{I\lambda(u)t'S_\emptyset}$. We need to find $t_m\rightarrow\infty$ such that $1_{I\lambda(u)I}1_{It't_mSp_{2n}}=1_{I\lambda(u)t't_mSp_{2n}}$. We can take $t_m\rightarrow\infty$ such that $t't_m$ is dominant and then this holds.

        Denote the sub algebra of $\Theta$ generated by $\theta_{\lambda}-q^\frac{l(\lambda^-)-l(\lambda^+)}{2}$ for $\lambda\in L_{max}$ by $\Theta_X$.

        We need to show that the mentioned elements generate the annihilator of $f_\emptyset$.

        Assume that $h\in H(G,I)$ satisfies $hf_\emptyset=0$. We can write $h=\sum_{w,\lambda} a_\lambda^wT_w\theta_\lambda$ for $w\in W$, $\lambda\in X_*(T)$ and coefficients $a_\lambda^w\in \C[q^{\pm 1}]$. After changing the coefficients by a factor of a power of $q$ we have $hf_\emptyset= \sum a_\lambda^w T_w1_{I\lambda(u)S_\emptyset}$.

        For fixed $\lambda$ and for every $w\in W$, $T_w1_{I\lambda(u)S_\emptyset}$ is supported on $IW\lambda(u)S_\emptyset$. We claim that for $\lambda,\lambda'$ such that $I\lambda(u)S_\emptyset\neq I\lambda(u)'S_\emptyset$ we have $IW\lambda(u)S_\emptyset\cap IW\lambda(u)'S_\emptyset=\emptyset$. To prove this it is enough to show that for every $\lambda''$ dominant enough if $I\lambda(u)\lambda(u)''Sp_{2n}\neq I\lambda(u)'\lambda(u)''Sp_{2n}$ then $IW\lambda(u)\lambda(u)''Sp_{2n}\cap IW\lambda(u)'\lambda(u)''Sp_{2n}=\emptyset$. 
        
        In any $I$ orbit in $IW\lambda(u)\lambda(u)''Sp_{2n}$ there is a unique anti-symmetric matrices whose non zero entries are equal to (up to a sign) $\lambda_1(u)\lambda''_1(u)\lambda_2(u)\lambda''_2(u),\lambda_3(u)\lambda''_3(u)\lambda_4(u)\lambda''_4(u),...$. 
        
        For $IW\lambda(u)\lambda(u)''Sp_{2n}$ the same is true with $\lambda_i$ replaced by $\lambda'_i$. By our assumption not all of these values are the same, and thus the $I$ orbits are different.

        Thus, in the expression $\sum_{w,\lambda} a_\lambda^w T_w1_{I\lambda(u)S_\emptyset}=0$ we can group the elements by the functions $1_{I\lambda(u)S_\emptyset}$ and for every fixed $\lambda$ we have $\sum_w a_\lambda^w T_w1_{I\lambda(u)S_\emptyset}=0$. 

        We claim that if $h'\in H_f$ and $h'1_{I\lambda(u)S_\emptyset}=0$ then $h'$ is in the ideal generated by $T_{s_{2i-1}}-q$. It is easy to see that $T_{s_{2i-1}}-q$ acts as zero on $1_{I\lambda(u)S_\emptyset}$. The dimension of $H_f1_{I\lambda(u)S_\emptyset}$ is equal to the number of $I$ orbits in $IW\lambda(u)S_\emptyset$. The number of such orbits is equal to the size of the coset space $W/\langle s_{2i-1}|i=1,...,2n\rangle$ and the result follows. 

        We proved that an element in the annihilator of $f_\emptyset$ modulo $\Theta_X$ can be written as sum of elements in $H_f(T_{s_{2i-1}}-q)\Theta$. To complete the proof we use Theorem \ref{Ber_presentation}.

        For $s=s_{2i-1}$ and $\alpha=\alpha_{2i-1}$ we have $$(T_s-q)\theta_\lambda=\theta_{s(\lambda)}T_s+(1-q)\frac{\theta_{s(\lambda)}-\theta_{\lambda}}{1-\theta_{-\alpha^\vee}}-q\theta_\lambda=\theta_{s(\lambda)}(T_s-q)+q\theta_{s(\lambda)}+(1-q)\frac{\theta_{s(\lambda)}-\theta_{\lambda}}{1-\theta_{-\alpha^\vee}}-q\theta_\lambda$$
        
        We have $q\theta_{s(\lambda)}+(1-q)\frac{\theta_{s(\lambda)}-\theta_{\lambda}}{1-\theta_{-\alpha^\vee}}-q\theta_\lambda=\frac{\theta_{s(\lambda)}-\theta_{\lambda}}{1-\theta_{-\alpha^\vee}}(q-q\theta_{-\alpha^\vee}+1-q)=\frac{\theta_{s(\lambda)}-\theta_{\lambda}}{1-\theta_{-\alpha^\vee}}(-q\theta_{-\alpha^\vee}+1)$. 

        Notice that $(-q\theta_{-\alpha^\vee}+1)\in \Theta_X$.
        
        We see that $H_f(T_{s_{2i-1}}-q)\Theta$ is in $H(G,I)\Theta_x+H(G,I)(T_{s_{2i-1}}-q)$ which completes the proof.
        
    \end{proof}

\end{document}
